\subsection{PASSAGE FROM LIE ALGEBRAS TO LIE GROUPS}

We shall denote the Hausdorff series by H(X, Y) (Chapter II, § 6, no. 4, Definition 1).

Lemma 2. Let \(\mathbf{L}\) be a complete normed Lie algebra over \(\mathbf{R}\) or \(\mathbf{C}\) . Let \(\mathbf{G}\) be the set of \(x \in \mathbf{L}\) such that \(\| x \| < \frac{1}{3} \log \frac{3}{2}\) . Let \(\theta\) be the mapping \(x \mapsto -x\) of \(\mathbf{G}\) into \(\mathbf{G}\) . Let \(\mathbf{H}\) be the restriction to \(\mathbf{G} \times \mathbf{G}\) of the Hausdorff function of \(\mathbf{L}\) (Chapter II, § 7, no. 2).

\begin{enumerate}
\def\labelenumi{(\roman{enumi})}
\item
  (G, 0, θ, H) is a Lie group germ.
\item
  Let \(\phi\) be the identity mapping of G into L. The differential of \(\phi\) at 0 is an isomorphism of the normable Lie algebra L(G) onto L.
\end{enumerate}
(i) follows from Chapter II, § 7, no. 2.

As \(\phi\) is a chart on G, the differential \(\psi\) of \(\phi\) at 0 is an isomorphism of normable spaces. On the other hand, the expansion as an integral series \(\mathrm{H} = \sum_{i,j\geqslant 0}\mathrm{H}_{ij}\) of the mapping H is such that \(\mathrm{H}_{11}(x,y) = \frac{1}{2} [x,y]\) . By § 3, Proposition 24, for all \(a,b\) in \(\mathrm{L}(\mathrm{G})\) ,

\[
\psi ([ a, b ]) = \mathrm{H} _ {1 1} (\psi (a), \psi (b)) - \mathrm{H} _ {1 1} (\psi (b), \psi (a)) = [ \psi (a), \psi (b) ]
\]

which proves (ii).

G is called the Lie group germ defined by L.

Suppose that K is ultrametric. Let p be the characteristic of the residue field of K. If \(p \neq 0\) , let \(\lambda = |p|^{1/(p-1)}\) ; if p = 0, let \(\lambda = 1\) .

Lemma 3. Let L be a complete normed Lie algebra over K. Let G be the set of \(x \in \mathbf{L}\) such that \(\| x \| < \lambda\) . Let H: G × G → G be the Hausdorff function of L (Chapter II, § 8, no. 3).

\begin{enumerate}
\def\labelenumi{(\roman{enumi})}
\item
  With the law of composition H, G is a Lie group in which 0 is identity element and -x the inverse of x for all \(x \in G\) .
\item
  Let \(\phi\) be the identity mapping of \(\mathbf{G}\) into \(\mathbf{L}\) . The differential of \(\phi\) at 0 is an isomorphism of the normable Lie algebra \(\mathrm{L}(\mathrm{G})\) onto \(\mathbf{L}\) .
\item
  For all \(\mu \in \mathbf{R}_+^*\) , let \(\mathrm{G}_{\mu}\) be the set of \(x \in \mathbf{L}\) such that \(\| x \| < \mu\) . Then the \(\mathrm{G}_{\mu}\) for \(\mu < \lambda\) , form a fundamental system of open and closed neighbourhoods of 0 and are subgroups of \(\mathbf{G}\) .
\end{enumerate}

Assertions (i) and (iii) follow from Chapter II, § 8, no. 3, Proposition 3 and (ii) can be proved as in Lemma 2.

G is called the Lie group defined by L.

THEOREM 2. Let \(\mathbf{L}\) be a complete normable Lie algebra. There exists a Lie group germ \(G\) such that \(\mathrm{L}(\mathbf{G})\) is isomorphic to \(\mathbf{L}\) . Two such Lie group germs are locally isomorphic.

The first assertion follows from Lemmas 2 and 3. The second assertion follows from Corollary 1 to Theorem 1 of no. 1.

COROLLARY 1. Let \(\mathbf{G}\) be a Lie group. There exists a neighbourhood of \(e\) which contains no finite subgroup distinct from \(\{e\}\) . If \(\mathbf{K} = \mathbf{R}\) or \(\mathbf{C}\) , there exists an open neighbourhood of \(e\) which contains no subgroup distinct from \(\{e\}\) .

We write \(\mathrm{L}(G) = \mathrm{L}\) . Choose a norm on \(\mathbf{L}\) defining the topology on \(\mathbf{L}\) and such that \(\| [x,y]\| \leqslant \| x\| \| y\|\) for all \(x,y\) in \(\mathbf{L}\) .

Suppose that \(\mathbf{K} = \mathbf{R}\) or \(\mathbf{C}\) . Let \(G'\) be the Lie group germ defined by L. There exist an open ball \(U'\) of centre 0 in \(G'\) and an isomorphism \(\phi\) of the Lie group germ \(U'\) onto an open neighbourhood \(U\) of \(e\) in G. Let \(V' = \frac{1}{2} U'\) , \(V = \phi(V')\) , H be a subgroup of G contained in V and \(h \in H\) . Let \(x = \phi^{-1}(h) \in V'\) . If \(x \neq 0\) , there exists an integer \(n > 0\) such that \(x, 2x, \ldots, nx\) are in \(V'\) , \((n + 1)x \in U'\) , \((n + 1)x \notin V'\) . Then \(h, h^2, \ldots, h^n\) are in V,

\[
h ^ {n + 1} \in \mathrm{U}, \quad h ^ {n + 1} \notin \mathrm{V},
\]

which is absurd. Hence \(\mathbf{H} = \{e\}\) .

Suppose that K is ultrametric. It suffices to prove the corollary when G is the Lie group associated with L. If \(g \in G\) , the powers of g evaluated in G are the elements of Zg evaluated in L. These are distinct if \(g \neq e\) . Hence G contains no finite subgroup distinct from \(\{e\}\) .

COROLLARY 2. Let \(k\) be a non-discrete closed subfield of \(\mathbf{K}\) , \(\mathbf{G}\) a Lie group over \(k\) and \(\mathbf{L} = \mathbf{L}(\mathbf{G})\) . Suppose that \(\mathbf{L}\) has a normable Lie \(\mathbf{K}\) -algebra structure \(\mathbf{L}'\) , compatible with the normable Lie \(k\) -algebra structure and invariant under the adjoint representation of \(\mathbf{G}\) . Then there exists on \(\mathbf{G}\) one and only one Lie \(\mathbf{K}\) -group structure compatible with the Lie \(k\) -group structure and for which the Lie algebra is \(\mathbf{L}'\) .

There exists a Lie group germ \(G_{1}\) over K such that \(\mathrm{L}(G_{1}) = \mathrm{L}'\) (Theorem 2). By Corollary 1 to Theorem 1 of no. 1, G and \(G_{1}\) , considered as Lie k-group germs, are locally isomorphic. Hence there exist an open neighbourhood \(G'\) of e in G and a Lie K-group germ structure on \(G'\) , with Lie algebra L, which is compatible with the Lie group germ structure over k. Let V be a symmetric open neighbourhood of e in G such that \(V^{2} \subset G'\) . Let \(g \in G\) . Then \(\phi = Int g\) is a k-isomorphism of a sufficiently small open Lie subgroup germ of \(G'\) onto an open Lie subgroup germ of \(G'\) ; and \(T_{e}(\phi)\) is K-linear and hence \(T_{x}(\phi)\) is K-linear for x sufficiently close to e; therefore the restriction of Int g to a sufficiently small open neighbourhood of e in V is K-analytic (Differentiable and Analytic Manifolds, R, 5.14.6). By §1, no. 9, Proposition 18, there exists on G an analytic K-manifold structure with which G is a Lie

K-group and V is an open sub-K-manifold of G. By translation, it is seen that the underlying k-manifold structure of G is the given structure. The Lie algebra of the Lie K-group G is the same as that of the open Lie K-subgroup germ V and hence is L'. Finally, the uniqueness stated in the corollary follows from § 3, no. 8, Proposition 32.

THEOREM 3. Let G be a Lie group germ and \(\mathfrak{h}\) a Lie subalgebra of \(\mathrm{L}(\mathrm{G})\) admitting a topological supplement. There exists a Lie subgroup germ H of G such that \(\mathrm{L}(\mathrm{H}) = \mathfrak{h}\) . If \(\mathrm{H}_1\) and \(\mathrm{H}_2\) are Lie subgroup germs of G such that

\[
\mathrm{L} (\mathrm{H} _ {1}) = \mathrm{L} (\mathrm{H} _ {2}) = \mathfrak {h},
\]

then \(\mathrm{H}_1\cap \mathrm{H}_2\) is open in \(\mathrm{H}_1\) and \(\mathrm{H}_2\)

There exists a Lie group germ H' with Lie algebra isomorphic to \(\mathfrak{h}\) (Theorem 2). Shrinking H' if necessary, it can be assumed that there exists a morphism \(\phi\) of H' into G such that \(\mathrm{L}(\phi)\) is an isomorphism of \(\mathrm{L}(\mathrm{H}')\) onto \(\mathfrak{h}\) (no. 1, Theorem 1). As \(\mathfrak{h}\) admits a topological supplement, \(\phi\) is an immersion at e. Hence, shrinking H' further, it can be assumed that \(\phi\) is an isomorphism of the manifold H' onto a submanifold of G. This proves the existence of H. The second assertion follows from the following proposition:

PROPOSITION 1. Let G be a Lie group germ and H and H' two Lie subgroup germs. In order that L(H) ⊃ L(H'), it is necessary and sufficient that H ∩ H' be open in H'. If H ∩ H' is open in H', then L(H') = L(H ∩ H') ⊂ L(H). Suppose that L(H) ⊃ L(H'). Let i, i' be the canonical injections of H, H' into G. By shrinking H' if necessary, it can be assumed that there exists a morphism ψ of H' into H such that L(ψ) is the canonical injection of L(H') into L(H) (no. 1, Theorem 1). Then L(i ∘ ψ) = L(i') and hence there exists a neighbourhood V of eH, in H' such that i ∘ ψ and i' coincide on V (Theorem 1). Therefore V ⊂ H, hence V ⊂ H ∩ H' and H ∩ H' is open in H' (§ 1, no. 10).

PROPOSITION 2. Let G be a Lie group over K, k a non-discrete closed subfield of K and H a Lie subgroup of the Lie k-group G. Suppose that L(H) is a vector sub-K-space of L(G) which admits a topological supplement. Then H is a Lie subgroup of the Lie K-group G.

There exists a Lie subgroup germ \(\mathrm{H}'\) of the Lie K-group G such that \(\mathrm{L}(\mathrm{H}') = \mathrm{L}(\mathrm{H})\) (Theorem 3). Consider G, H, H' as Lie \(k\) -group germs; Theorem 3 then proves that \(\mathrm{H} \cap \mathrm{H}'\) is open in H and H'. Hence there exists an open neighbourhood U of \(e\) in G such that \(\mathrm{U} \cap \mathrm{H}\) is a submanifold of G over K. Therefore, H is a Lie subgroup of the Lie K-group G ( \(\S 1\) , no. 3, Proposition 6).

